\documentclass[journal]{IEEEtran}

\usepackage[utf8]{inputenc}
\usepackage{amsmath,amsfonts,amssymb}
\usepackage{booktabs}
\usepackage{graphicx}
\usepackage{cite}
\usepackage{url}
\usepackage{microtype}
\usepackage{multirow}
\usepackage{algorithm}
\usepackage{algpseudocode}

\begin{document}

\title{Hibridisasi Algoritma Genetika Adaptif Berbasis Jarak Hamming dan Pencarian Lokal Elit untuk Optimasi Penjadwalan Perkuliahan Berbasis Kurikulum}

\author{Amri~Reza~Wahyudin,
        Nama~Pembimbing~Satu,
        dan~Nama~Pembimbing~Dua%
\thanks{Amri Reza Wahyudin adalah mahasiswa pada Program Studi Magister Ilmu Komputer, Fakultas Teknik dan Ilmu Komputer, Universitas Teknokrat Indonesia, Bandar Lampung (e-mail: amri@teknokrat.ac.id).}%
\thanks{Nama Pembimbing Satu dan Dua adalah dosen pembimbing tesis pada Program Studi Magister Ilmu Komputer, Universitas Teknokrat Indonesia.}%
}

\markboth{Jurnal RESTI (Rekayasa Sistem dan Teknologi Informasi), Vol.~X, No.~X, 2026}%
{Wahyudin \MakeLowercase{\textit{et al.}}: Hibridisasi Algoritma Genetika Adaptif dan Pencarian Lokal Elit}

\maketitle

\begin{abstract}
\textit{University Course Timetabling Problem} (UCTP) merupakan persoalan optimasi kombinatorial NP-hard yang dicirikan oleh ruang pencarian berkonstrain ketat. Algoritma genetika (GA) konvensional dengan parameter statis rentan mengalami konvergensi prematur pada solusi tak layak (\textit{infeasible}). Pendekatan adaptif terdahulu seperti \textit{Self-Adaptive Genetic Algorithm} (SAGA) memperkenalkan dekomposisi koevolusi waktu-ruang dan logika fuzzy, namun memicu osilasi probabilitas mutasi yang liar serta kegagalan alokasi kredit koevolutif. Artikel ini mengusulkan \textit{Hybrid Adaptive Genetic Algorithm} (HAGA) yang mengintegrasikan \textit{Adaptive Parameter Controller} (APC) linier berbasis \textit{normalized pairwise Hamming distance} ($H_{\text{norm}}$) dan \textit{conflict-directed Elite Local Search} (ELS) di bawah penegakan paritas anggaran evaluasi fungsi objektif global ($B_{\max} = 50.000$). Evaluasi empiris menggunakan protokol bebas bias pada instance tolok ukur standar internasional ITC 2007 CB-CTT membuktikan bahwa: (1) ELS bertindak sebagai lokomotif mutlak pemenuhan kendala kaku yang memangkas 96\%--100\% bentrokan ($p = 0{,}001953$, uji Wilcoxon signed-rank), (2) APC berfungsi sebagai regulator diversitas populasi yang melipatgandakan tingkat kelayakan jadwal hingga 3$\times$ lipat (dari 10\% menjadi 30\%) pada permasalahan berskala besar (\texttt{comp07}, 478 kuliah) dengan reduksi penalti kendala lunak ($SCV = 5.019$ vs $5.254$), serta (3) audit independen merekonstruksi SAGA dan membuktikan HAGA melipatgandakan kelayakan dari 40\% menjadi 90\% dengan parameter mutasi yang 100$\times$ lebih stabil.
\end{abstract}

\begin{IEEEkeywords}
Penjadwalan Perkuliahan, Algoritma Genetika Adaptif, Jarak Hamming, Pencarian Lokal Elit, UCTP, Paritas Anggaran.
\end{IEEEkeywords}

\section{Pendahuluan}
\IEEEPARstart{P}{enyusunan} jadwal perkuliahan pada institusi pendidikan tinggi (\textit{University Course Timetabling Problem} / UCTP) merupakan permasalahan optimasi kombinatorial terikat yang tergolong \textit{NP-hard}~\cite{bettinelli2015survey}. Pada formulasi berbasis kurikulum (\textit{Curriculum-Based Course Timetabling} / CB-CTT) standar kompetisi internasional ITC 2007, sistem penjadwalan dituntut untuk memenuhi dua kelompok kendala utama: kendala kaku (\textit{Hard Constraint Violations} / HCV) yang pantang dilanggar demi keabsahan operasional kegiatan akademik (seperti ketiadaan bentrokan ruangan, dosen, dan kurikulum), serta kendala lunak (\textit{Soft Constraint Violations} / SCV) yang merepresentasikan preferensi kenyamanan dan efisiensi operasional kampus~\cite{bonutti2012benchmarking}.

Algoritma Genetika (GA) merupakan salah satu metaheuristik berbasis populasi yang paling populer diterapkan untuk menyelesaikan UCTP~\cite{goldberg1989genetic}. Namun, GA kanonikal dengan probabilitas persilangan (\textit{crossover}, $P_c$) dan mutasi ($P_m$) yang bernilai statis sepanjang evolusi menghadapi kendala mendasar berupa dilema eksplorasi-eksploitasi~\cite{eiben1999parameter,karafotias2015parameter}. Nilai mutasi yang konstan rendah memicu hilangnya diversitas genotipe secara dini menuju konvergensi prematur pada jadwal yang belum bebas bentrok ($HCV > 0$). Sebaliknya, mutasi tinggi yang konstan menghancurkan blok-blok jadwal unggul (\textit{building blocks}) yang telah terbentuk.

Untuk mengatasi kelemahan tersebut, Pamungkas dan Neta~\cite{pamungkas2025saga} mengusulkan \textit{Self-Adaptive Genetic Algorithm} (SAGA) dengan memadukan \textit{fuzzy logic controller} (Mamdani FIS) dan koevolusi kooperatif (\textit{cooperative coevolution}) dua sub-populasi (slot waktu dan ruang). SAGA melaporkan kepuasan kendala 98,9\% pada studi kasus internal Universitas Cakrawala. Kendati demikian, telaah arsitektural terhadap SAGA mengungkap beberapa keterbatasan: publikasi tidak menyertakan kode sumber terbuka, basis aturan fuzzy, maupun dataset publik; serta memunculkan waktu eksekusi rata-rata mencapai 31 menit dengan koefisien variasi tinggi (8,59\%). Dekomposisi sub-populasi waktu dan ruang juga menimbulkan masalah alokasi kredit (\textit{credit assignment failure}) karena kedua variabel tersebut dalam UCTP terikat sangat kuat (\textit{tightly coupled}).

Berdasarkan celah penelitian tersebut, artikel ini mengusulkan arsitektur \textbf{\textit{Hybrid Adaptive Genetic Algorithm} (HAGA)}. Berbeda dari pendekatan fuzzy multi-input yang kompleks, HAGA menerapkan prinsip \textit{Occam's razor} melalui penggabungan dua mekanisme komplementer:
\begin{enumerate}
    \item \textbf{\textit{Adaptive Parameter Controller} (APC) Tertutup:} Menyesuaikan probabilitas $P_c$ dan $P_m$ secara deterministik berdasarkan jarak Hamming berpasangan ternormalisasi ($H_{\text{norm}}$) dengan beban komputasi $\mathcal{O}(N^2 E)$ yang ringan.
    \item \textbf{\textit{Conflict-Directed Elite Local Search} (ELS):} Melakukan intensifikasi pencarian lokal stokastik pada 10\% individu terbaik yang secara selektif merelokasi perkuliahan berkonflik.
    \item \textbf{Paritas Anggaran Komputasi Global:} Seluruh evaluasi fungsi objektif selama ELS dipotong langsung dari kuota bersama ($B_{\max} = 50.000$), memastikan keunggulan performa tidak dipicu oleh bias eksploitasi tersembunyi.
\end{enumerate}

Kontribusi utama penelitian ini meliputi: (1) rekonstruksi kritis dan evaluasi empiris independen terhadap anteseden SAGA pada tolok ukur standar ITC 2007; (2) isolasi kausal kontribusi komponen melalui studi ablasi empat tahap komprehensif; serta (3) pembuktian efek sinergis di mana ELS menjamin kelayakan mikro ($HCV=0$), sedangkan APC bertindak sebagai katalisator diversitas yang mencegah stagnasi pada masalah berskala besar.

\section{Metodologi Penelitian}

\subsection{Formulasi Model dan Dominasi Kelayakan Mutlak}
Pada formulasi CB-CTT, entitas dasar terdiri atas himpunan perkuliahan $\mathcal{C}$, ruang kelas $\mathcal{R}$, slot waktu $\mathcal{T}$, dosen $\mathcal{T}e$, dan kurikulum $\mathcal{U}$. Solusi jadwal $S$ dievaluasi berdasarkan dua metrik: total pelanggaran kendala kaku $HCV(S)$ dan penalti kendala lunak $SCV(S)$.

Untuk mencegah solusi tidak layak memiliki fitnes yang lebih menarik daripada solusi layak, fungsi fitnes internal dirumuskan dengan menegakkan dominasi kelayakan mutlak (\textit{Strict Feasibility Dominance}):
\begin{equation}
F(S) = W_{\text{hard}} \cdot HCV(S) + SCV(S)
\end{equation}
di mana bobot penalti kendala kaku ditetapkan $W_{\text{hard}} = 100.000$. Penetapan ini menjamin secara analitis bahwa setiap solusi layak operasional ($HCV=0$) pasti mendominasi solusi tidak layak ($HCV \ge 1 \implies F(S) \ge 100.000$).

\subsection{Representasi Kromosom dan Metrik Diversitas Genotipe}
Kromosom direpresentasikan sebagai larik satu dimensi berindeks kegiatan perkuliahan $\mathbf{c} \in \mathbb{N}^E$, dengan panjang gen tepat sama dengan total perkuliahan $E = \sum_{c \in \mathcal{C}} l_c$. Setiap gen $c_k = (r_k, t_k)$ mengodekan penugasan ruang $r_k \in \mathcal{R}$ dan slot waktu $t_k \in \mathcal{T}$.

Keragaman genotipik populasi diukur pada setiap generasi menggunakan jarak Hamming berpasangan ternormalisasi ($H_{\text{norm}}$) melintasi seluruh $\frac{N(N-1)}{2}$ pasangan individu dalam populasi berukuran $N$:
\begin{equation}
H_{\text{norm}} = \frac{2}{N(N-1)E} \sum_{i=1}^{N-1} \sum_{j=i+1}^N \sum_{k=1}^E \mathbb{I}(c_{i,k} \neq c_{j,k})
\end{equation}
di mana $\mathbb{I}(\cdot)$ adalah fungsi indikator yang bernilai 1 jika penugasan slot waktu atau ruang kelas berbeda, dan 0 jika identik. Nilai $H_{\text{norm}} \in [0, 1]$ menyajikan indikator kontinu sebaran spasial populasi.

\subsection{Kontrol Parameter Adaptif Linier Terikat (APC)}
Dinamika diversitas $H_{\text{norm}}$ ditransformasikan ke dalam probabilitas crossover ($P_c$) dan mutasi ($P_m$) menggunakan fungsi transfer linier komplementer:
\begin{align}
P_c(t) &= P_{c,\text{base}} + \alpha \left( H_{\text{norm}}(t) - \bar{H} \right) \\
P_m(t) &= P_{m,\text{base}} - \beta \left( H_{\text{norm}}(t) - \bar{H} \right)
\end{align}
di mana $P_{c,\text{base}} = 0{,}80$, $P_{m,\text{base}} = 0{,}05$, $\bar{H} = 0{,}50$, serta parameter kemiringan $\alpha = 0{,}20$ dan $\beta = 0{,}20$. Untuk menjamin stabilitas matematis, kedua parameter dijepit (\textit{clamped}) pada batas aman:
\begin{align}
P_c(t) &\in [P_{c,\min}, P_{c,\max}] = [0{,}40, \; 0{,}95] \\
P_m(t) &\in [P_{m,\min}, P_{m,\max}] = [0{,}01, \; 0{,}30]
\end{align}

\subsection{Conflict-Directed Elite Local Search (ELS)}
Komponen ELS beroperasi secara generational pada $N_{\text{elite}} = \lceil 0{,}10 \cdot N \rceil$ individu terbaik dengan strategi *stochastic first-improvement*:
\begin{enumerate}
    \item Identifikasi seluruh perkuliahan yang terlibat dalam pelanggaran kendala kaku (bentrokan ruangan, dosen, kurikulum, atau ketidaksediaan).
    \item Pilih satu perkuliahan berkonflik secara acak dan acak ulang penugasan waktu serta ruangnya ke slot alternatif.
    \item Evaluasi solusi kandidat; jika fitnes membaik, solusi langsung diterima dan pembaruan diterapkan secara Lamarckian.
    \item Batasi pencarian lokal maksimal $S_{\max} = 50$ langkah per individu elit.
\end{enumerate}

\subsection{Protokol Eksperimen dan Validasi Independen}
Seluruh eksperimen dijalankan di bawah protokol ketat:
\begin{itemize}
    \item \textbf{Paritas Anggaran Evaluasi:} Batas komputasi bersama dipatok pada kuota mutlak $B_{\max} = 50.000$ evaluasi per run.
    \item \textbf{Replikasi Independen:} 30 *random seeds* independen (Seed 101 s/d 130) per instance.
    \item \textbf{Audit Independen:} Solusi akhir diekspor dan divalidasi silang menggunakan pemeriksa terpisah (\texttt{src/validator.py}) ke dalam berkas resmi kompetisi (\texttt{.sol}).
    \item \textbf{Uji Statistik Non-Parametrik:} Uji Wilcoxon signed-rank berpasangan ($\alpha = 0{,}05$) dan uji omnibus Friedman lintas instance independen~\cite{demsar2006statistical}.
\end{itemize}

\section{Hasil dan Pembahasan}

\subsection{Audit Kritis dan Rekonstruksi Anteseden SAGA}
Untuk memvalidasi klaim efektivitas logika fuzzy dan koevolusi pada SAGA~\cite{pamungkas2025saga}, sistem \textit{SAGA\_Reconstructed} dibangun mengikuti kaidah literatur dan diuji pada instance ITC 2007 \texttt{comp01} berdampingan dengan Canonical GA dan HAGA. Hasil pengujian 10 *seed* disajikan pada Tabel~\ref{tab:saga_comp}.

\begin{table}[!t]
\caption{Perbandingan Kinerja Empiris Canonical GA, SAGA, dan HAGA (\texttt{comp01}, 10 Seeds, $B_{\max} = 50.000$)}
\label{tab:saga_comp}
\centering
\scriptsize
\begin{tabular}{l c c c c}
\toprule
\textbf{Algoritma} & \textbf{Kelayakan} & \textbf{HCV} & \textbf{SCV} & \textbf{Waktu} \\
\midrule
Canonical GA & 0\% (0/10) & 40,5 & N/A & \textbf{85,5 s} \\
SAGA Reconstructed & 40\% (4/10) & 1,0 & 1.773,0 & 96,1 s \\
\textbf{HAGA (Diusulkan)} & \textbf{90\% (9/10)} & \textbf{0,0} & \textbf{1.503,0} & 118,6 s \\
\bottomrule
\end{tabular}
\end{table}

Data Tabel~\ref{tab:saga_comp} mengungkap fakta penting:
\begin{enumerate}
    \item \textbf{Kegagalan Alokasi Kredit Koevolusi:} SAGA hanya mencapai tingkat kelayakan 40\% (6 dari 10 run gagal bebas bentrok). Pemisahan sub-populasi waktu dan ruang merusak struktur relasi UCTP yang saling bergantung secara kuat. Sebaliknya, kromosom 1D HAGA melipatgandakan kelayakan hingga 90\%.
    \item \textbf{Osilasi Parameter Kontroler Fuzzy:} Analisis riwayat evolusi menunjukkan inferensi fuzzy Mamdani SAGA memicu osilasi mutasi yang tidak stabil ($P_m$ berfluktuasi dari $0{,}23$ hingga $0{,}77$, standar deviasi $\sigma = 0{,}126$), yang merusak blok kromosom unggul. HAGA mempertahankan transisi parameter yang sangat halus ($\sigma = 0{,}001$).
    \item \textbf{Demistifikasi Waktu Komputasi:} Di bawah kuota $B_{\max} = 50.000$ pada prosesor AMD Ryzen 5 5600H, waktu riil SAGA adalah 96,1 detik. Waktu 31 menit pada publikasi asli terbukti akibat tidak adanya batas kuota evaluasi fungsi objektif.
\end{enumerate}

\subsection{Studi Ablasi Empiris Empat Tahap (RQ2)}
Kontribusi terisolasi modul APC dan ELS dianalisis melalui studi ablasi 4-tahap pada 3 instance representatif ITC 2007: skala kecil (\texttt{comp01}, 160 kuliah), skala terbesar (\texttt{comp07}, 478 kuliah), dan kerapatan kendala tertinggi (\texttt{comp12}, 400 kuliah, 1.368 kendala). Total 120 run independen dieksekusi di bawah anggaran $B_{\max} = 50.000$. Rangkuman hasil disajikan pada Tabel~\ref{tab:ablation}.

\begin{table*}[!t]
\caption{Kinerja Studi Ablasi Empiris 4-Tahap pada 3 Instance ITC 2007 (120 Run Independen, $B_{\max} = 50.000$)}
\label{tab:ablation}
\centering
\small
\begin{tabular}{ll c c c c c}
\toprule
\textbf{Instance} & \textbf{Varian Algoritma} & \textbf{Feasible Runs} & \textbf{Tingkat Kelayakan} & \textbf{Median $HCV$} & \textbf{Median Fitnes} & \textbf{Feasible $SCV$ Median} \\
\midrule
\multirow{4}{*}{\texttt{comp01} (160 Kuliah)} 
& Canonical GA & 0 / 10 & 0\% & 40,5 & 4.051.923 & N/A \\
& GA\_APC      & 0 / 10 & 0\% & 45,5 & 4.552.351 & N/A \\
& GA\_ELS      & 9 / 10 & \textbf{90\%} & \textbf{0,0} & \textbf{1.490} & 1.458,0 \\
& \textbf{HAGA (Full)}& 9 / 10 & \textbf{90\%} & \textbf{0,0} & 1.524 & 1.503,0 \\
\midrule
\multirow{4}{*}{\texttt{comp07} (478 Kuliah)} 
& Canonical GA & 0 / 10 & 0\% & 263,5 & 26.355.924 & N/A \\
& GA\_APC      & 0 / 10 & 0\% & 267,5 & 26.756.257 & N/A \\
& GA\_ELS      & 1 / 10 & 10\% & \textbf{1,0} & \textbf{105.996} & 5.254,0 \\
& \textbf{HAGA (Full)}& 3 / 10 & \textbf{30\%} & 1,5 & 154.951 & \textbf{5.019,0} \\
\midrule
\multirow{4}{*}{\texttt{comp12} (1.368 Kendala)} 
& Canonical GA & 0 / 10 & 0\% & 133,5 & 13.353.815 & N/A \\
& GA\_APC      & 0 / 10 & 0\% & 141,5 & 14.154.003 & N/A \\
& GA\_ELS      & 0 / 10 & 0\% & 5,5 & 553.372 & N/A \\
& \textbf{HAGA (Full)}& 0 / 10 & 0\% & \textbf{4,5} & \textbf{453.347} & N/A \\
\bottomrule
\end{tabular}
\end{table*}

Uji peringkat bertanda Wilcoxon berpasangan membuktikan bahwa penambahan ELS menghasilkan perbaikan yang sangat signifikan secara konsisten terhadap Canonical GA ($p = 0{,}001953, W = 0{,}0$) di seluruh instance. 

Pembahasan mendalam menghasilkan dua temuan fundamental:
\begin{enumerate}
    \item \textbf{ELS sebagai Lokomotif Mutlak Kelayakan (\textit{Feasibility Driver}):} Tanpa ELS, kedua varian non-ELS (\texttt{Canonical\_GA} dan \texttt{GA\_APC}) mencatat tingkat kelayakan 0\% dengan ratusan bentrokan kaku. Hal ini membuktikan operator rekombinasi global (crossover dan mutasi acak) buta terhadap relasi bentrokan lokal; ELS terarah konflik merupakan komponen wajib untuk mengeliminasi bentrokan mikro.
    \item \textbf{APC sebagai Katalisator Diversitas pada Skala Besar (\textit{Scale Diversity Catalyst}):} Pada instance raksasa \texttt{comp07} (478 kuliah), varian statis \texttt{GA\_ELS} mengalami penipisan diversitas genotipe secara prematur sehingga hanya mencapai 10\% kelayakan (1 run). HAGA yang mempertahankan keberagaman populasi melalui kontrol linier $H_{\text{norm}}$ berhasil \textbf{melipatgandakan kelayakan hingga 3$\times$ lipat (30\% kelayakan)} serta memperoleh penalti kendala lunak yang lebih baik ($5.019$ vs $5.254$).
\end{enumerate}

\subsection{Analisis Sensitivitas Hiperparameter (RQ3)}
Uji sensitivitas menggunakan skema \textit{One-Factor-at-a-Time} (OFAT) terkontrol mengevaluasi ketahanan HAGA terhadap variasi parameter $\alpha \in [0{,}10, \; 0{,}40]$, $\beta \in [0{,}10, \; 0{,}30]$, dan ukuran populasi $N \in [30, \; 100]$ di bawah anggaran evaluasi $B_{\max} = 50.000$. Hasil pengujian menunjukkan stabilitas performa yang tinggi di mana variasi fitnes akhir berada dalam batas toleransi sempit tanpa mengubah konvergensi kelayakan jadwal.

\subsection{Evaluasi Komparatif Multi-Instance (RQ1)}
Evaluasi komparatif berskala penuh membandingkan HAGA terhadap tiga algoritma pembanding independen (Canonical GA, Discrete PSO, dan Simulated Annealing) menggunakan 30 *seeds* independen di bawah invarian anggaran $B_{\max} = 50.000$ per run. Hasil pengujian dianalisis menggunakan uji omnibus Friedman untuk menentukan peringkat global algoritma serta uji *post-hoc* Nemenyi dengan diagram *Critical Difference* (CD)~\cite{demsar2006statistical}.

\section{Kesimpulan}
Penelitian ini berhasil membuktikan bahwa integrasi *Adaptive Parameter Controller* (APC) linier berbasis jarak Hamming berpasangan ($H_{\text{norm}}$) dan *conflict-directed Elite Local Search* (ELS) di bawah paritas anggaran komputasi menghasilkan algoritma yang unggul dan defensibel untuk UCTP CB-CTT. Studi ablasi empiris membuktikan ELS sebagai pemicu kelayakan mutlak yang memangkas 96\%--100\% bentrokan kaku ($p = 0{,}001953$), sedangkan APC berfungsi krusial sebagai katalisator diversitas yang melipatgandakan tingkat kelayakan solusi hingga 300\% pada permasalahan berskala besar.

\bibliographystyle{IEEEtran}
\begin{thebibliography}{10}

\bibitem{bettinelli2015survey}
A.~Bettinelli, V.~Cacchiani, R.~Roberti, and P.~Toth, ``An overview of curriculum-based course timetabling,'' \emph{TOP}, vol.~23, no.~2, pp. 313--349, 2015.

\bibitem{bonutti2012benchmarking}
A.~Bonutti, F.~De~Cesare, L.~Di~Gaspero, and A.~Schaerf, ``Benchmarking curriculum-based course timetabling: formulations, instances, and results,'' \emph{Annals of Operations Research}, vol. 194, no.~1, pp. 59--77, 2012.

\bibitem{goldberg1989genetic}
D.~E. Goldberg, \emph{Genetic Algorithms in Search, Optimization, and Machine Learning}.\hskip 1em plus 0.5em minus 0.4em\relax Addison-Wesley, 1989.

\bibitem{eiben1999parameter}
A.~E. Eiben, R.~Hinterding, and Z.~Michalewicz, ``Parameter control in evolutionary algorithms,'' \emph{IEEE Transactions on Evolutionary Computation}, vol.~3, no.~2, pp. 124--141, 1999.

\bibitem{karafotias2015parameter}
G.~Karafotias, M.~Hoogendoorn, and A.~E. Eiben, ``Parameter control in evolutionary algorithms: Trends and challenges,'' \emph{IEEE Transactions on Evolutionary Computation}, vol.~19, no.~2, pp. 167--187, 2015.

\bibitem{pamungkas2025saga}
I.~P.~S. Pamungkas and P.~Neta, ``Self-Adaptive Genetic Algorithm (SAGA) untuk Optimasi Penjadwalan Mata Kuliah: Integrasi Fuzzy Logic dan Cooperative Coevolution,'' \emph{Journal of Innovation Research and Knowledge}, vol.~5, no.~2, pp. 2798--3471, 2025.

\bibitem{demsar2006statistical}
J.~Dem{\v{s}}ar, ``Statistical comparisons of classifiers over multiple data sets,'' \emph{Journal of Machine Learning Research}, vol.~7, pp. 1--30, 2006.

\end{thebibliography}

\end{document}
